\section{Z-Flagged Double Checks}\label{sec:z_flagged_double_checks}

So how can the intended fault distance be restored?
Note the distance-3 double check uses 6 ancillas
to interact with the 7 data qubits
so there is necessarily a pair of CXs
from the same ancilla to different data qubits (e to 5, and e to 6 via f).
One may therefore think using 7 ancillas will restore fault distance;
however,
none of the X hook errors in \cref{fig:all_2_fault_configurations}
are due to this CX pair.
Varying ancilla count will not restore fault distance.

A solution which works is to replace the double check
with two short or long single checks,
sketched in \cref{fig:short_single_check_sketch,fig:long_single_check_sketch}
respectively.
We find
using the analysis of the previous section
that these $\overline{H_+}$-measurement methods
do not have degraded fault distance,
at least for code distance 3;
specifically, we analyse the distance-3 short single check
(Crumble link \href{https://algassert.com/crumble\#circuit=Q(0,0)0;Q(0,1)1;Q(0,2)2;Q(0,3)3;Q(0,4)4;Q(1,0)5;Q(1,1)6;Q(1,3)7;Q(1,4)8;Q(2,0)9;Q(2,1)10;Q(2,2)11;Q(2,3)12;Q(2,4)13;Q(3,0)14;Q(3,1)15;Q(3,2)16;Q(3,3)17;Q(3,4)18;Q(4,0)19;Q(4,1)20;Q(4,2)21;Q(4,3)22;Q(4,4)23;MPP_Y6*Y10*Y15*Y11*Y17*Y12*Y7;POLYGON(0,0,1,0.25)7_6_11_12;POLYGON(0,1,0,0.25)10_12_17;POLYGON(1,0,0,0.25)6_11_15;TICK;MPP_X6*X11*X12*X7;TICK;MPP_X6*X10*X15*X11;TICK;MPP_X11*X10*X17*X12;TICK;MPP_Z6*Z11*Z12*Z7;TICK;MPP_Z6*Z10*Z15*Z11;TICK;MPP_Z11*Z10*Z17*Z12;TICK;R_23;RX_22;TICK;CX_22_23;R_21_18;TICK;CX_22_21_23_18;R_13_20;TICK;CX_18_13_21_20;MX_22_23;R_8_19;TICK;CX_13_8_20_19;R_4_14;TICK;CX_8_4_19_14;MX_20;R_9_3;TICK;CX_14_9_4_3;MX_8_19;R_5_2;TICK;CX_3_2_9_5;MX_4;R_0_1_16;TICK;CX_5_0_2_1_21_16;TICK;CX_1_0;MX_5_2_21;S_6_10_12;S_DAG_7_11_17_15;TICK;CX_1_6_3_7_13_12_18_17_9_10_14_15_16_11;M_0;DT(0,0,0)rec[-1];TICK;MX_1_9_14_16_3_13_18;S_11_7_17_15;S_DAG_10_6_12;DT(1,1,1)rec[-1]_rec[-2]_rec[-3]_rec[-4]_rec[-5]_rec[-6]_rec[-7]_rec[-9]_rec[-10]_rec[-11]_rec[-12]_rec[-13]_rec[-14]_rec[-15]_rec[-16]_rec[-17]_rec[-24];TICK;MPP_X6*X11*X12*X7;DT(1,1,2)rec[-1]_rec[-24];TICK;MPP_X6*X10*X15*X11;DT(1,1,3)rec[-1]_rec[-24];TICK;MPP_X11*X10*X17*X12;DT(2,2,4)rec[-1]_rec[-11]_rec[-24];TICK;MPP_Z6*Z11*Z12*Z7;DT(1,1,5)rec[-1]_rec[-24];TICK;MPP_Z6*Z10*Z15*Z11;DT(1,1,6)rec[-1]_rec[-24];TICK;MPP_Z11*Z10*Z17*Z12;DT(2,2,7)rec[-1]_rec[-14]_rec[-24];TICK;MPP_Y6*Y10*Y15*Y11*Y17*Y12*Y7;OI(0)rec[-1]_rec[-8]_rec[-9]_rec[-10]_rec[-11]_rec[-12]_rec[-13]_rec[-14]_rec[-16]_rec[-17]_rec[-18]_rec[-19]_rec[-20]_rec[-21]_rec[-22]_rec[-23]_rec[-24]}{here})
from the predecessor of MSC, Itogawa et al.~\cite{Itogawa2025},
and a minimally modified long single-check version
(Crumble link \href{https://algassert.com/crumble\#circuit=Q(0,0)0;Q(0,1)1;Q(0,3)2;Q(0,4)3;Q(1,0)4;Q(1,1)5;Q(1,2)6;Q(1,3)7;Q(1,4)8;Q(2,0)9;Q(2,1)10;Q(2,2)11;Q(2,3)12;Q(2,4)13;Q(3,0)14;Q(3,1)15;Q(3,2)16;Q(3,3)17;Q(3,4)18;MPP_Y1*Y5*Y10*Y6*Y12*Y7*Y2;POLYGON(0,0,1,0.25)2_1_6_7;POLYGON(0,1,0,0.25)5_7_12;POLYGON(1,0,0,0.25)1_6_10;DT(0,1,0)rec[-1];TICK;MPP_X1*X6*X7*X2;TICK;MPP_X1*X5*X10*X6;TICK;MPP_X6*X5*X12*X7;TICK;MPP_Z1*Z6*Z7*Z2;TICK;MPP_Z1*Z5*Z10*Z6;TICK;MPP_Z6*Z5*Z12*Z7;TICK;R_17;RX_16;TICK;CX_16_17;R_18_15;TICK;CX_17_18_16_15;R_14_13_11;TICK;CX_15_14_18_13_16_11;R_8_9;TICK;CX_14_9_13_8;R_4_3;TICK;CX_8_3_9_4;R_0;TICK;CX_4_0;S_1_5_7;S_DAG_2_6_12_10;TICK;CX_8_7_13_12_4_5_9_10_11_6_0_1_3_2;TICK;CX_4_0;S_6_2_12_10;S_DAG_5_1_7;TICK;CX_8_3_9_4;M_0;DT(0,0,1)rec[-1];TICK;CX_14_9_13_8;M_4_3;DT(1,0,2)rec[-2];DT(0,4,2)rec[-1];TICK;CX_15_14_18_13_16_11;M_9_8;DT(2,0,3)rec[-2];DT(1,4,3)rec[-1];TICK;CX_17_18_16_15;M_14_11_13;DT(3,0,4)rec[-3];DT(2,2,4)rec[-2];DT(2,4,4)rec[-1];TICK;CX_16_17;M_15_18;DT(3,1,5)rec[-2];DT(3,4,5)rec[-1];TICK;M_17;MX_16;DT(3,3,6)rec[-2];DT(0,1,6)rec[-1]_rec[-19];TICK;MPP_X1*X6*X7*X2;DT(0,1,7)rec[-1]_rec[-19];TICK;MPP_X1*X5*X10*X6;DT(0,1,8)rec[-1]_rec[-19];TICK;MPP_X6*X5*X12*X7;DT(1,2,9)rec[-1]_rec[-19];TICK;MPP_Z1*Z6*Z7*Z2;DT(0,1,10)rec[-1]_rec[-19];TICK;MPP_Z1*Z5*Z10*Z6;DT(0,1,11)rec[-1]_rec[-19];TICK;MPP_Z6*Z5*Z12*Z7;DT(1,2,12)rec[-1]_rec[-19];TICK;MPP_Y1*Y5*Y10*Y6*Y12*Y7*Y2;OI(0)rec[-1]_rec[-8]}{here}).
Sahay et al.~\cite[\S B.3.c]{Sahay2026} use the latter technique,
exploiting 2D-nonlocal connectivity
to pipeline the (un)encoding of the GHZ states
for its two consecutive long single checks.
Under 2D-local connectivity this pipelining is unavailable,
so it is desirable to keep the double check
especially since it uses nearly half the spacetime volume
per logical measurement compared to the long single check.

We therefore offer an alternative solution,
\cref{fig:d3a6f2,fig:d5a19f13},
which keeps the double check but adds Z flags
to detect all low-fault-count malignant configurations
(see \cref{fig:d3a6f2_conventional_diagram} for
the conventional circuit diagram of \cref{fig:d3a6f2}).
Our Z flags are 2D local in that
	each interacts only with a pair of neighbouring qubits,
but they are not 2D \emph{nearest-neighbour} since
	the existing interactions become next-nearest-neighbour.
We therefore do not intend for them to be physically implemented verbatim,
but rather we treat them as a proof of concept.
In \cref{sec:pre_escape_performance} we show
Monte Carlo simulation results with and without these flags,
and in \cref{sec:flag_design_process}
we briefly comment on our flag design process.

\begin{figure}
	\centering
	\begin{tikzpicture}
	\node[anchor=south west, inner sep=0] (image)
		{\includegraphics[width=0.9\textwidth]{figures/d3a6f2_T.png}};
	\begin{scope}[x={(image.south east)}, y={(image.north west)}]
		\foreach \x/\y in {0.511/0.766, 0.767/0.766, 0.00/0.001, 0.256/0.001}
		\draw[draw=zblue, line width=1.5pt]
			(\x,\y) rectangle ++(0.234,0.233);
	\end{scope}
	\end{tikzpicture}
	\caption{A Z-flagged version of the distance-3 double-check circuit
	that detects the four fault configurations
	shown in \cref{fig:all_2_fault_configurations}.
	The two added flag qubits are those prepared in $\ket{0}$ in Tick 2,
	and the four additional CX layers are boxed in blue.
	Red, green, and blue shapes represent
	$\X$, $\H_\pm :=\frac{1}{\sqrt{2}}(\X \pm \Y)$, and $\Z$
	\emph{detecting regions} (defined in \cref{sec:reducing_code_conversion_to_codespace_projection}),
	respectively.
	Click \href{https://algassert.com/crumble\#circuit=Q(0,0)0;Q(0,1)1;Q(1,0)2;Q(1,1)3;Q(1.5,0.5)4;Q(2,0)5;Q(2,1)6;Q(2,2)7;Q(2,3)8;Q(2.5,0.5)9;Q(3,0)10;Q(3,1)11;Q(3,2)12;Q(4,0)13;Q(4,1)14;POLYGON(0,0,1,0.25)10_13_12_6;POLYGON(0,1,0,0.25)8_12_6_3;POLYGON(1,0,0,0.25)3_6_10_0;TICK;MPP_Y0*Y10*Y13*Y12*Y8*Y6*Y3;TICK;MPP_X10*X0*X3*X6;TICK;MPP_X3*X6*X12*X8;TICK;MPP_X13*X10*X6*X12;TICK;MPP_Z10*Z0*Z3*Z6;TICK;MPP_Z8*Z12*Z6*Z3;TICK;MPP_Z13*Z10*Z6*Z12;TICK;RX_14_11_5_2_7_1;S_3_10_12;S_DAG_0_6_8_13;TICK;CX_1_0_11_10_7_8_5_6_2_3_14_13;R_4_9;TICK;CX_3_4_6_9;TICK;CX_6_4_11_9;TICK;CX_3_1_6_7_11_14;TICK;CX_6_3_11_12;TICK;CX_6_11;TICK;MX_6;DT(0,0,0)rec[-1]_rec[-8];TICK;RX_6;TICK;CX_6_11;TICK;CX_6_3_11_12;TICK;CX_3_1_6_7_11_14;TICK;CX_6_4_11_9;TICK;CX_3_4_6_9;TICK;CX_1_0_11_10_7_8_14_13_5_6_2_3;M_4_9;DT(1.5,0.5,1)rec[-2];DT(2.5,0.5,1)rec[-1];TICK;MX_14_11_5_2_7_1;S_6_13_8_0;S_DAG_3_10_12;DT(4,1,2)rec[-6];DT(3,1,2)rec[-5];DT(2,1,2)rec[-4]_rec[-9];DT(1,0,2)rec[-3];DT(2,2,2)rec[-2];DT(0,1,2)rec[-1];TICK;MPP_X10*X0*X3*X6;DT(3,0,3)rec[-1]_rec[-10]_rec[-16];TICK;MPP_X8*X12*X6*X3;DT(1,1,4)rec[-1]_rec[-11]_rec[-16];TICK;MPP_X13*X10*X6*X12;DT(4,0,5)rec[-1]_rec[-12]_rec[-16];TICK;MPP_Z10*Z0*Z3*Z6;DT(3,0,6)rec[-1]_rec[-16];TICK;MPP_Z8*Z12*Z6*Z3;DT(2,3,7)rec[-1]_rec[-16];TICK;MPP_Z13*Z10*Z6*Z12;DT(4,0,8)rec[-1]_rec[-16];TICK;MPP_Y0*Y10*Y13*Y12*Y8*Y6*Y3;OI(0)rec[-1]_rec[-8]_rec[-9]_rec[-12]_rec[-13]}{here}
	for the Crumble link.}
	\label{fig:d3a6f2}
\end{figure}

\begin{figure}
	\centering
	\begin{tikzpicture}
	\node[anchor=south west, inner sep=0] (image)
		{\includegraphics[width=0.97\textwidth]{figures/d5a19f13_T.png}};
	\begin{scope}[x={(image.south east)}, y={(image.north west)}]
		\foreach \x/\y in {0.512/0.814, 0.256/0.001}
		\draw[draw=zblue, line width=1.5pt]
			(\x,\y) rectangle ++(0.234,0.185);
	\end{scope}
	\end{tikzpicture}
	\caption{A Z-flagged version of the distance-5 double-check circuit
	that has fault distance 5.
	The 13 added flag qubits are those prepared in $\ket{0}$,
	and the two additional circuit layers are boxed in blue.
	Click \href{https://algassert.com/crumble\#circuit=Q(0,0)0;Q(0,1)1;Q(1,0)2;Q(1,1)3;Q(1.5,0.5)4;Q(2,0)5;Q(2,1)6;Q(2,2)7;Q(2,3)8;Q(2,4)9;Q(2.5,0.5)10;Q(2.5,1.5)11;Q(2.5,2.5)12;Q(3,0)13;Q(3,1)14;Q(3,2)15;Q(3,3)16;Q(3,4)17;Q(3.5,0.5)18;Q(4,0)19;Q(4,1)20;Q(4,2)21;Q(4,3)22;Q(4,4)23;Q(4,5)24;Q(4,6)25;Q(4.5,0.5)26;Q(4.5,1.5)27;Q(4.5,2.5)28;Q(4.5,3.5)29;Q(4.5,4.5)30;Q(5,0)31;Q(5,1)32;Q(5,2)33;Q(5,3)34;Q(5,4)35;Q(5,5)36;Q(5,6)37;Q(5.5,0.5)38;Q(5.5,1.5)39;Q(6,0)40;Q(6,1)41;Q(6,2)42;Q(6,3)43;Q(6.5,0.5)44;Q(7,0)45;Q(7,1)46;Q(7,2)47;Q(7,3)48;Q(8,0)49;Q(8,1)50;POLYGON(0,0,1,0.25)13_19_32_21_15_6;POLYGON(0,0,1,0.25)41_45_49_47;POLYGON(0,0,1,0.25)34_43_36_23;POLYGON(0,1,0,0.25)8_15_6_3;POLYGON(0,1,0,0.25)32_41_47_43_34_21;POLYGON(0,1,0,0.25)17_23_36_25;POLYGON(1,0,0,0.25)3_6_13_0;POLYGON(1,0,0,0.25)19_45_41_32;POLYGON(1,0,0,0.25)15_21_34_23_17_8;TICK;MPP_X0*X13*X6*X3_X19*X45*X41*X32_X15*X21*X34*X23*X17*X8;TICK;MPP_X32*X41*X47*X43*X34*X21_X3*X6*X15*X8_X17*X23*X36*X25;TICK;MPP_X13*X19*X32*X21*X15*X6_X34*X43*X36*X23_X41*X45*X49*X47;TICK;MPP_Z0*Z13*Z6*Z3_Z19*Z45*Z41*Z32_Z15*Z21*Z34*Z23*Z17*Z8;TICK;MPP_Z32*Z41*Z47*Z43*Z34*Z21_Z3*Z6*Z15*Z8_Z17*Z23*Z36*Z25;TICK;MPP_Z13*Z19*Z32*Z21*Z15*Z6_Z34*Z43*Z36*Z23_Z41*Z45*Z49*Z47;TICK;MPP_Y3*Y6*Y0*Y13*Y19*Y32*Y45*Y41*Y47*Y43*Y34*Y36*Y25*Y23*Y17*Y8*Y15*Y21*Y49;TICK;RX_31_40_2_5_7_9_22_24_48_1_14_20_46_42_16_33_35_37_50;S_3_13_32_15_17_36_34_47_45;S_DAG_19_21_23_25_41_43_49_6_8_0;TICK;CX_1_0_2_3_5_6_14_13_20_19_31_32_40_41_46_45_42_43_48_47_33_21_22_34_35_23_24_36_37_25_7_15_16_8_9_17_50_49;R_11_27_30_29_12_18;TICK;CX_16_12_35_29_36_30_15_11_33_27_14_18;R_39_28_44_38_4_10;TICK;CX_16_17_36_37_14_11_32_27_33_28_35_30_34_29_15_12_42_39_46_44_41_38_3_4_6_10;TICK;CX_35_36_15_16_14_10_42_47_33_39_32_38_34_28_46_50_41_44_3_1_6_4;R_26;TICK;CX_34_35_14_15_33_42_20_18_41_46_32_26_6_3;TICK;CX_14_6_33_34_20_26_32_41;TICK;CX_20_14_32_33;TICK;CX_32_20;TICK;MX_32;DT(1,1,0)rec[-1]_rec[-2];TICK;RX_32;TICK;CX_32_20;TICK;CX_20_14_32_33;TICK;CX_14_6_33_34_20_26_32_41;TICK;CX_34_35_14_15_33_42_20_18_41_46_32_26_6_3;TICK;CX_35_36_15_16_14_10_42_47_33_39_32_38_34_28_46_50_41_44_3_1_6_4;M_26;DT(4.5,0.5,1)rec[-1];TICK;CX_16_17_36_37_14_11_32_27_33_28_35_30_34_29_15_12_42_39_46_44_41_38_3_4_6_10;TICK;CX_16_12_35_29_36_30_15_11_33_27_14_18;M_39_28_44_38_4_10;DT(5.5,1.5,2)rec[-6];DT(4.5,2.5,2)rec[-5];DT(6.5,0.5,2)rec[-4];DT(5.5,0.5,2)rec[-3];DT(1.5,0.5,2)rec[-2];DT(2.5,0.5,2)rec[-1];TICK;CX_1_0_2_3_5_6_14_13_20_19_31_32_40_41_46_45_42_43_48_47_33_21_22_34_35_23_24_36_37_25_16_8_9_17_7_15_50_49;M_11_27_30_29_12_18;DT(2.5,1.5,3)rec[-6];DT(4.5,1.5,3)rec[-5];DT(4.5,4.5,3)rec[-4];DT(4.5,3.5,3)rec[-3];DT(2.5,2.5,3)rec[-2];DT(3.5,0.5,3)rec[-1];TICK;MX_22_24_7_9_2_5_31_40_48_1_14_20_46_42_16_33_35_37_50;S_19_21_23_25_41_43_6_8_0_49;S_DAG_3_13_32_15_17_36_34_47_45;DT(4,3,4)rec[-19];DT(4,5,4)rec[-18];DT(2,2,4)rec[-17];DT(2,4,4)rec[-16];DT(1,0,4)rec[-15];DT(2,0,4)rec[-14];DT(5,1,4)rec[-13]_rec[-33];DT(6,0,4)rec[-12];DT(7,3,4)rec[-11];DT(0,1,4)rec[-10];DT(3,1,4)rec[-9];DT(4,1,4)rec[-8];DT(7,1,4)rec[-7];DT(6,2,4)rec[-6];DT(3,3,4)rec[-5];DT(5,2,4)rec[-4];DT(5,4,4)rec[-3];DT(5,6,4)rec[-2];DT(8,1,4)rec[-1];TICK;MPP_X0*X13*X6*X3_X19*X45*X41*X32_X15*X21*X34*X23*X17*X8;DT(0,0,5)rec[-3]_rec[-55];DT(4,0,5)rec[-2]_rec[-36]_rec[-54];DT(3,2,5)rec[-1]_rec[-53];TICK;MPP_X32*X41*X47*X43*X34*X21_X3*X6*X15*X8_X17*X23*X36*X25;DT(5,1,6)rec[-3]_rec[-39]_rec[-55];DT(1,1,6)rec[-2]_rec[-54];DT(3,4,6)rec[-1]_rec[-53];TICK;MPP_X13*X19*X32*X21*X15*X6_X34*X43*X36*X23_X45*X49*X47*X41;DT(3,0,7)rec[-3]_rec[-42]_rec[-55];DT(5,3,7)rec[-2]_rec[-54];DT(6,1,7)rec[-1]_rec[-53];TICK;MPP_Z0*Z13*Z6*Z3_Z19*Z45*Z41*Z32_Z15*Z21*Z34*Z23*Z17*Z8;DT(0,0,8)rec[-3]_rec[-55];DT(4,0,8)rec[-2]_rec[-54];DT(3,2,8)rec[-1]_rec[-53];TICK;MPP_Z32*Z41*Z47*Z43*Z34*Z21_Z3*Z6*Z15*Z8_Z17*Z23*Z36*Z25;DT(5,1,9)rec[-3]_rec[-55];DT(1,1,9)rec[-2]_rec[-54]}{here}
	for the Crumble link.}
	\label{fig:d5a19f13}
\end{figure}
\begin{SCfigure}
	\includegraphics[width=0.65\linewidth]{figures/d3_pre_escape_ler.pdf}
	\caption{The logical error rate per kept shot of
	distance-3 cultivation before the escape stage.
	Blue (orange) data correspond to cultivating S (T) states.
	Shading shows the region within a factor of 1000 of the maximum likelihood.}
	\label{fig:d3_pre_escape_ler}
\end{SCfigure}

\begin{SCfigure}
	\includegraphics[width=0.47\linewidth]{figures/improvement_and_cost_ratios.pdf}
	\caption{The pre-escape performance of distance-5 T-state cultivation,
	with and without Z flags.
	Turquoise:
	the logical error rate ratio (original / flagged);
	sand:
	the expected attempts ratio (flagged / original);
	shading shows 95\% Bayesian credible intervals.
	This figure and \cref{fig:d5_pre_escape_ler} are derived from the same simulation data.}
	\label{fig:improvement_and_cost_ratios}
\end{SCfigure}

\subsection{Pre-Escape Performance}\label{sec:pre_escape_performance}
To see how our Z-flagged MSC performs in practice,
we use the fast universal simulator SymFT~\cite{Fang2026a}
to simulate original and Z-flagged MSC
up to and excluding the escape stage.
\Cref{fig:d3_pre_escape_ler} shows that our Z flags recover
the intended fault distance of MSC-3,
thus pinning the final double check as
the sole culprit degrading its fault distance.
For MSC-3 in practice we can see that these Z flags are overkill,
since the malignant 2-fault configurations they detect
barely outweigh the higher-fault-count configurations they introduce.

For MSC-5 however \cref{fig:d5_pre_escape_ler,fig:improvement_and_cost_ratios} show they are practical:
in addition to the numbers stated in the abstract,
at noise level $p =\num{5e-4}$
the Z flags improve the pre-escape logical error rate by 10.9$\times$,
costing a 1.17$\times$ increase in attempts per kept shot.
Note the simulated circuits have an intended fault distance of 4
as \cref{sec:magic_state_cultivation} explains,
so the true fault distances are 3 without and 4 with flags.
However,
a weighted-least-squares fit gives log--log slopes of
3.8 without and 5.1 with flags,
indicating the simulated noise levels do not yet reach the asymptotic scaling regime.
Our accompanying code~\cite{Chan2026c_quantum_bibstyle}
also includes an unsimulated circuit for Z-flagged MSC-5 with an intended fault distance of 5
i.e.\ with four instead of three stabiliser rounds before its distance-5 double check.

The analysis from Wan and Zapirain~\cite[Appendix~H.3--4]{Wan2026a}
more rigorously verifies that with our Z flags:
MSC-3,
three-round MSC-5,
and four-round MSC-5
indeed have fault distances 3, 4, and 5, respectively.
In the third case,
our Z flags restore the pre-escape logical error rate
from $\upTheta(p^3)$ to $\upTheta(p^5)$.
Pre-escape logical error rates are relevant to full MSC
if the escape-stage decoder confidence threshold is set high:
Chase and Labib~\cite[\S 4.2.2]{Chase2026} showed that in this regime,
the T--S discrepancy persists after the escape stage,
thus most of the end-to-end logical error
comes from the injection and cultivation stages
rather than the purely Clifford escape stage.

\subsection{Flag Design Process}\label{sec:flag_design_process}
For distance 3:
we added the minimum number of Z flags,
preferring short-range interactions,
that together detects each X fault in \cref{fig:all_2_fault_configurations}.
For distance 5:
there are too many ($3 +601$) malignant configurations of \num{<5} faults
to individually design flags for.
Instead,
we continued the trend from distance 3,
adding a Z flag between each of the 18 consecutive qubit pairs
in the tree along which the GHZ state is (un)encoded;
we verified that the resulting double check has fault distance 5.
We then removed the flag at each of the 5 leaves of the tree
as we found this did not reduce fault distance.
Finally,
we minimised the lifetime, hence the idling error, of each flag qubit.
\Cref{alg:probability_trivial_syndrome,prop:resulting_fidelity}
were crucial as they enabled the fast verification of whether
a candidate Z-flagged double check had the intended fault distance.